Dans un spectromètre de masse, des ions une fois positifs ($q = e = \qty{1.602e-19}{C}$), initialement presque au repos, sont accélérés par une différence de potentiel de \qty{\v{dV}}{kV}. Ils entrent ensuite perpendiculairement dans un champ magnétique uniforme de \qty{\v{B}}{T}, décrivent un demi-cercle et frappent un détecteur placé dans le plan d'entrée. Le faisceau contient deux isotopes de masses \qty{\v{m1}}{u} et \qty{\fpeval{\v{m1}+\v{dm}}}{u} ($\qty{1}{u} = \qty{1.661e-27}{kg}$).

\begin{enumerate}[label=\alph*)]
    \item Déterminez la vitesse de l'isotope le plus léger à l'entrée du champ magnétique.
    \item Déterminez le rayon de la trajectoire de chacun des deux isotopes.
    \item Déterminez la distance qui sépare les deux points d'impact sur le détecteur.
\end{enumerate}
